\uuid{1d5g}
\exo7id{3563}
\titre{exo7 3563}
\auteur{quercia}
\organisation{exo7}
\datecreate{2010-03-10}
\isIndication{false}
\isCorrection{true}
\chapitre{Réduction d'endomorphisme, polynôme annulateur}
\sousChapitre{Polynôme annulateur}
\cpge{mpsi-mp2i.polynomes-fractions-rationnelles.c, pcsi.polynomes.c, ptsi.polynomes.c, mpsi-mp2i.changements-bases.c, bcpst1.matrices, bcpst1.polynomes-reels.b, mp-mpi.reduction-endomorphismes.g, pc.complements-espaces-vectoriels.d, psi.complements-espaces-vectoriels.d, bcpst2.polynomes.b, bcpst2.valeurs-propres-vecteurs-propres.a}
\module{Algèbre}
\niveau{L2}
\difficulte{}

\contenu{
\question{
Caractériser les polynômes $P$ tels que~: $\forall\ A\in\mathcal{M}_n(\C)$, $(P(A) = 0)  \Rightarrow  (\mathrm{tr}(A)\in\Z)$.
}
\reponse{
Aucun polynôme constant ne convient. Si $P$ est non constant et $\alpha$
est une racine de~$P$ alors en considérant $A=\alpha I_n$ on obtient une
première condition nécessaire~: $n\alpha \in\Z$. Si $P$ a une autre racine $\beta$
alors en prenant $A = \mathrm{diag}(\alpha,\dots,\alpha,\beta)$ on obtient
une deuxième condition nécessaire~: $\beta-\alpha\in\Z$. Ainsi les polynômes~$P$
cherchés ont la propriété suivante~: {\it $\deg(P)\ge 1$ et il existe $u\in\Z$ tel que
toutes les racines de~$P$ sont congrues à $u/n$ modulo~$1$.}
Cette condition est clairement suffisante.
}
}
